\documentclass[../main.tex]{subfiles}

\begin{document}

\subsection{Why the Sequence Persuades}
\label{subsec:story-structure}

The sequence persuades because it prevents the reader from jumping too quickly
to a positive conclusion about inclusion. The chapter begins with structural
conditions, moves through readiness and access, tests whether those conditions
produce real digital usage, and only then asks whether that usage translates
into resilience. The same story can still be summarized once as
\textit{Context \(\rightarrow\) Readiness \(\rightarrow\) Access \(\rightarrow\)
Usage \(\rightarrow\) Resilience}, but the force of the sequence lies in what
each transition corrects.

The key corrective turn happens after Page 3. Banking Inclusion reaches
\(74.93\%\), and Financial Institution Account ownership reaches \(74.05\%\),
which could easily support a success narrative. Page 4 interrupts that reading
by showing Agricultural Sales at \(73.42\%\) cash-only and Urban Merchant
Payments at \(65.79\%\) versus \(50.30\%\) in rural areas. Page 5 then
interrupts any remaining optimism by showing that the Poorest 20\% still rely
far more on Family \& Friends at 86 than on Formal Savings at 30. In other
words, the chapter persuades because it steadily narrows the distance between
headline access and lived financial security.

\noindent\textit{See Figures~\ref{fig:account-kpi-cards},
\ref{fig:national-inflows}, and
\ref{fig:page5-coping-income-group}.}

\end{document}
